% 型7：答案添削・採点基準型
% 実際の答案を載せ，採点基準表と講評を添える．
% 記述式で「何を書けば点になるか」を伝える用途に向く．
\documentclass[lualatex,paper=b5j,fontsize=10pt,twoside]{jlreq}

\usepackage{surikumi-mondaishu}

\MathMagazineSetup{
  feature={},
  series={答案研究／数学 II},
  series-position=left,
  pattern=parquet,
  title={不等式の証明で何が点になるか},
  author={入試基礎講座},
  author-font=mincho,
  author-position=right,
  title-fill=black,
  deck={結論が合っていても書き落とすと減点される箇所がある．採点基準を先に見てから答案を読むこと．}
}

\begin{document}

\MakeMathMagazineTitle

\begin{mmproblem}[1]{}
$a>0$，$b>0$ のとき，
\[
  \dfrac{a+b}{2}\geq\sqrt{ab}
\]
が成り立つことを示せ．また，等号が成り立つ条件を求めよ．
\end{mmproblem}

\begin{msrubric}
  差 $\dfrac{a+b}{2}-\sqrt{ab}$ を考える方針を示している & 2 \\
  $(\sqrt{a}-\sqrt{b})^2$ の形に変形している & 3 \\
  $a>0$，$b>0$ より $\sqrt{a}$，$\sqrt{b}$ が実数であることに触れている & 1 \\
  平方が $0$ 以上であることから不等式を結論している & 2 \\
  等号成立条件を $a=b$ と正しく求めている & 2 \\
\end{msrubric}

\MMSubheading{答案例}

\begin{mmsolution}[答案]
$a>0$，$b>0$ であるから $\sqrt{a}$ と $\sqrt{b}$ はともに実数である．
このとき
\begin{align*}
  \dfrac{a+b}{2}-\sqrt{ab}
    &= \dfrac{a-2\sqrt{ab}+b}{2}\\
    &= \dfrac{(\sqrt{a}-\sqrt{b})^2}{2}
\end{align*}
となる．実数の平方は $0$ 以上であるから
$(\sqrt{a}-\sqrt{b})^2\geq 0$ であり，したがって
\[
  \dfrac{a+b}{2}-\sqrt{ab}\geq 0
\]
が成り立つ．すなわち $\dfrac{a+b}{2}\geq\sqrt{ab}$ である．

等号が成り立つのは $(\sqrt{a}-\sqrt{b})^2=0$ のとき，
すなわち $\sqrt{a}=\sqrt{b}$ のときである．
$a>0$，$b>0$ より，これは
\[
  a=b
\]
と同値である．
\end{mmsolution}

\MSComment{差をとって平方の形に持ち込む流れは書けている．
$\sqrt{a}$，$\sqrt{b}$ が実数であることへの一言があるため，
平方が $0$ 以上という根拠が宙に浮いていない．満点である．}

\MMSubheading{減点される答案}

\begin{msmistake}[減点例]
両辺を 2 乗すると $\dfrac{(a+b)^2}{4}\geq ab$ となる．
展開して $a^2+2ab+b^2\geq 4ab$ より $(a-b)^2\geq 0$ であり，
これは常に成り立つ．よって示された．
\end{msmistake}

\MSWhereWrong{示したい不等式を最初に真と仮定して変形している．
このままでは「$(a-b)^2\geq 0$ ならばもとの不等式が成り立つ」ことを
示していない．同値変形であることを断るか，$(a-b)^2\geq 0$ から
出発して逆向きに書き直す必要がある．}

\MSComment{両辺が非負であることを述べたうえで
「2 乗しても大小関係は変わらない」と明示すれば同値変形になり，
この方針でも満点になる．方針が悪いのではなく，論理の向きの断りが
抜けている．}

\end{document}
